\documentclass[10pt,a4paper]{amsart}
\usepackage[margin=19mm]{geometry}
\usepackage{amsmath,amssymb,mathtools}
\usepackage[scr=rsfs,cal=euler]{mathalfa}
\usepackage{xfrac}
\usepackage[unicode,hidelinks,bookmarks=false]{hyperref}
\newcommand{\ie}{i.e.\ }
\newcommand{\cf}{cf.\ }
\newcommand{\resp}{respectively\ }
\newcommand{\Subsection}{\subsection}
\providecommand{\nobreakdash}{\nobreak\hbox{-}}
\setlength{\emergencystretch}{2em}
\multlinegap0pt
\usepackage{amssymb}
\usepackage{float}
\usepackage{tikz}
\newtheorem{remark}{Remark}
\newtheorem{theorem}{Theorem}
\newcommand{\C}{\mathbb{C}}
\newcommand{\M}{\mathbb{M}}
\newcommand{\inpr}[2]{\langle{#1,#2}\rangle}
\title{A family of matrix flows converging to normal matrices}
\author{Masaki Izumi}
\date{Source: arXiv:2602.11837v1 (CC BY 4.0)}
\begin{document}
\maketitle

\noindent\emph{Source and licence.} A family of matrix flows converging to normal matrices. Version arXiv:2602.11837v1 (CC BY 4.0), distributed by arXiv under the Creative Commons Attribution 4.0 International licence, http://creativecommons.org/licenses/by/4.0/, as recorded in the arXiv licence metadata for this exact version and shown on the abstract page https://arxiv.org/abs/2602.11837v1. Full licence notice: \texttt{licenses/arxiv-2602.11837-CC-BY-4.0.txt}.\par\smallskip

\noindent\textit{Editorial scope.} Counted unit: the article's theorem holomorphic (source0.tex lines 963--977). The article's complete original proof of it is delivered with it (source0.tex lines 979--1086, one \texttt{proof} environment of 108 lines). No mathematics is written, repaired or re-derived: the statement and the proof are the article's own lines, in the article's own notation, and no retained line is substituted or re-flowed.  Retained uncounted as the source-local context the proof needs: lines 533--535; lines 646--676; lines 851--856.  Nothing was omitted from any retained span, so the package carries no omission record.  No retained line contains a citation, so the wrapper carries no bibliography and no bibliography entry.  No retained line contains a figure environment, an image-inclusion command or drawing code, so no image is delivered and the package declares no asset dependency.  Editorial formatting only, in addition: an AMS \texttt{amsart} preamble that restates the article's own declarations and macros used by the retained lines, namely the environments \texttt{remark}, \texttt{theorem} on separate counters, together with the standard packages those lines need.  The retained environments print in this document as Remark 1, Theorem 1 in order of appearance, whereas the article prints the same items with the numbering of its own full document.  The complete native source of the article as distributed in the arXiv e-print for this exact version is bundled unchanged as \texttt{sources/194497/source0.tex}.
\begin{equation}\label{upptriflow}
\frac{dY(t)}{dt}=[D(Y(t))+2P(Y(t)),\Lambda+Y(t)],\quad Y(0)=T-\Lambda,
\end{equation}

\begin{remark}\label{diagonalizable} 
Under the above assumption, the following hold: 
\begin{itemize}
\item[(1)] If $T$ as above is diagonalizable, then $T_{ij}=0$ for all $i<j$ with $\lambda_i=\lambda_j$.   
	Indeed, if there exist two indices $i<j$ with $\lambda_i=\lambda_j$ and $T_{ij}\neq 0$, 
	we can express $T-\lambda_iI$ as the block matrix 
	\[T-\lambda_iI=\left(
	\begin{array}{ccc}
		A &B &*  \\
		O &N &*  \\
		O &O &C 
	\end{array}
	\right),
	\]
	so that $A$ and $C$ are invertible matrices and $N$ is a non-zero nilpotent matrix. 
	Since \[(T-\lambda_iI)^2=\left(
	\begin{array}{ccc}
		A^2 &AB+ BN&*  \\
		O &N^2 &*  \\
		O &O &C^2 
	\end{array}
	\right),
	\]
	we get $\ker (T-\lambda_i I)^2\varsupsetneqq\ker (T-\lambda_i I)$, and thus $T$ is not diagonalizable. 
\item[(2)] If $\mu_1=0$, and $Y(t)=(y(t)_{ij})$ is the unique solution of Eq.(\ref{upptriflow}), 
we have $y_{ij}(t)=0$ for all $1\leq i<j\leq m^g_1$ and $t\geq 0$. 
Indeed, it follows from the uniqueness of the solution since we can show that Eq.(\ref{upptriflow}) with 
this additional condition has a solution. 
In the rest of this section, we always assume $y_{ij}=0$ for $1\leq i<j\leq m^g$. 
\end{itemize}
\end{remark}	

\begin{align}\label{criterion}
\lefteqn{[\Lambda,\varphi(\Lambda+Y)]_{ij}} \\
	 &=(\lambda_i-\lambda_j)\left(  \frac{\varphi^{[1]}_1(|\lambda_i|,|\lambda_j|)}{|\lambda_i|+|\lambda_j|}\overline{\lambda_i}
	 -\frac{\varphi^{[1]}_2(|\lambda_i|,|\lambda_j|)}{|\lambda_i|+|\lambda_j|}\overline{\lambda_j}\right)y_{ij}
	 +o(\|[\Lambda,Y]\|_2),\quad (Y\to 0),\nonumber
\end{align}

\begin{theorem}\label{holomorphic} Let the notation be as above and assume that $\varphi$ satisfies the condition (C1) 
with $a=0$ and $\varphi_i(0)=0$. 
We further assume that $\varphi_i(\sqrt{x})$, $i=1,2$, 
extend to $\psi_i\in C(R_{b^2,\varphi})$ for some $\varepsilon>0$ such that they are holomorphic 
on the interior of $R_{b^2,\varepsilon}$.  
Furthermore, we assume that either of the following 2 conditions: 
\begin{itemize}
\item[$(1)$] The following integral converges for $i=1,2$: 
\[\int_{-\varepsilon}^{\varepsilon}\frac{|\psi_i(it)|}{t^2}dt<\infty.\] 
\item[$(2)$] $\varphi_2=0$, $\varphi_1'(x)/x$ is bounded on $(0,b]$, and 
\[\int_{-\varepsilon}^{\varepsilon}\frac{|\psi_1(it)|}{|t|}dt<\infty.\] 
\end{itemize}
Then the flow $\{F^\varphi_t(T)\}_{t\geq 0}$ converges to a normal matrix 
as $t$ tends to $\infty$ for any $T\in \M_n[0,b]$.  
\end{theorem}

\begin{proof} We verify Eq.(\ref{criterion}). 
Assume $1\leq i<j\leq n$ and $\lambda_i\neq \lambda_j$. 
Note that thanks to the integrability assumption and $\psi_i(0)=0$, the holomorphic functional calculus 
	\[\varphi(\Lambda+Y)=\frac{1}{2\pi i}\int_{\partial R_{b^2,\varepsilon}}
	\left(\frac{\psi_1(z)}{zI-|\Lambda+Y|^2}-\frac{\psi_2(z)}{zI-|\Lambda^*+Y^*|^2}\right)
	dz\]
makes sense. 
Thus 
\begin{align*}[\Lambda,P(Y)]_{ij}
	& =\frac{1}{2\pi i}\int_{\partial R_{b^2,\varepsilon}} \inpr{[\Lambda,\frac{1}{\zeta I-|\Lambda+Y|^2}]e_j}{e_j}\psi_1(\zeta)d\zeta \\
	&-\frac{1}{2\pi i}\int_{\partial R_{b^2,\varepsilon}} \inpr{[\Lambda,\frac{1}{\zeta I-|\Lambda^*+Y^*|^2}]e_j}{e_j}\psi_2(\zeta)d\zeta.
\end{align*}

The first integral is 
\begin{align*}
	&\frac{1}{2\pi i}\int_{\partial R_{b^2,\varepsilon}} \inpr{
		\frac{1}{\zeta I-|\Lambda+Y|^2}
		[\Lambda,|\Lambda+Y|^2] 
		\frac{1}{\zeta I-|\Lambda+Y|^2} e_j}{e_i}
	\psi_1(\zeta)d\zeta \\
	&=\frac{1}{2\pi i}\int_{\partial R_{b^2,\varepsilon}} \inpr{
		\frac{1}{\zeta I-|\Lambda+Y|^2}
		(\Lambda^*[\Lambda,Y]+[\Lambda,Y^*]\Lambda) 
		\frac{1}{\zeta I-|\Lambda+Y|^2} e_j}{e_i}
	\psi_1(\zeta)d\zeta \\
	&+\frac{1}{2\pi i}\int_{\partial R_{b^2,\varepsilon}} \inpr{
		\frac{1}{\zeta I-|\Lambda+Y|^2}
		(Y^*[\Lambda,Y]+[\Lambda,Y^*]Y) 
		\frac{1}{\zeta I-|\Lambda+Y|^2} e_j}{e_i}
	\psi_1(\zeta)d\zeta, 
\end{align*}
which can be computed from 
\begin{align*}
	&\frac{1}{2\pi i}\int_{\partial R_{b^2,\varepsilon}} \inpr{
		\frac{1}{\zeta I-|\Lambda+Y|^2}
		E_{kl} 
		\frac{1}{\zeta I-|\Lambda+Y|^2} e_j}{e_i}	\psi_1(\zeta)d\zeta\\ \nonumber
	&=\frac{1}{2\pi i}\int_{\partial R_{b^2,\varepsilon}} \inpr{\frac{1}{\zeta I-|\Lambda+Y|^2} e_k}{e_i}\inpr{\frac{1}{\zeta I-|\Lambda+Y|^2} e_j}{e_l}	\psi_1(\zeta)d\zeta. \nonumber
\end{align*}
Thus all we have to verify is  
\begin{align}\label{limit} 
	&\lim_{Y\to 0}
	\frac{1}{2\pi i}\int_{\partial R_{b^2,\varepsilon}} \inpr{\frac{1}{\zeta I-|\Lambda+Y|^2} e_k}{e_i}\inpr{\frac{1}{\zeta I-|\Lambda+Y|^2} e_j}{e_l}	\psi_1(\zeta)d\zeta\\ \nonumber
	&=\delta_{k,i}\delta_{l,j}\frac{1}{2\pi i}\int_{\partial R_{b^2,\varepsilon}}\frac{\psi_1(\zeta)}{(\zeta-|\lambda_i|^2)(\zeta-|\lambda_j|^2)}d\zeta,
\end{align}
and its counterpart for $|\Lambda^*+Y^*|$ as the rest of the computation is routine work. 
Note that we have 
\[\frac{1}{2\pi i}\int_{\partial R_{b^2,\varepsilon}}\frac{\psi_1(\zeta)}{(\zeta-|\lambda_i|^2)(\zeta-|\lambda_j|^2)}d\zeta
=\psi^{[1]}_1(|\lambda_i|^2,|\lambda_j|^2)
=\frac{\varphi_1^{[1]}(|\lambda_i|,|\lambda_j|)}{|\lambda_i|+|\lambda_j|}.\]




In the case (1), the Lebesgue convergence theorem verifies Eq.(\ref{limit}). 

Assume (2) now. 
Remark \ref{diagonalizable},(2) shows $y(t)_{pq}=0$ for all $1\leq p<q\leq m^g_1$, $t\geq 0$. 
Thus it suffices to verify  Eq.(\ref{limit}) under the additional condition $y_{pq}=0$ 
for all $1\leq p<q\leq m^g_1$, which we assume in the following argument.  
Then $\ker|\Lambda+Y(t)|=\mathrm{span}\{e_p\}_{p=1}^{m^g_1}$, which implies $\varphi(\Lambda+Y)_{ij}=0$ 
for $i\leq m^g_1$. 
We assume $m^g_1< i$ now. 
Since the rank of $\Lambda+Y$ is $n-m^g_1$, we have $s_p(\Lambda+Y)>0$ for $1\leq p\leq n-m^a_1$, 
and $s_p(\Lambda+Y)=0$ for $p\geq n-m^g_1+1$. 
The continuity of the singular numbers implies $s_q(\Lambda+Y)\to |\lambda_{n+1-q}|$ as $Y$ tends to 0. 
Thus there exist $\delta>0$ and $0<c_1<c_2<b$ such that for all $Y$ with $\|Y\|_2<\delta$, we have 
$s_p(\Lambda+Y)>c_2$ for $p\leq n-m^g_1$ and $s_p(\Lambda+Y)<c_1$ for $p>n-m^g_1$. 
We assume $\|Y\|_2<\delta$ in the following argument. 

Let $Q(Y)=\chi_{[0,c_1]}(|\Lambda+Y|)$, where $\chi_{[0,c_1]}$ is the characteristic function of $[0,c_1]$. 
Since we can replace $\chi_{[0,c_1]}$ with a Lipschitz function with the same values on $[0,c_1]$ and $[c_2,b]$, 
there exists $L>0$ satisfying $\|Q(Y)-Q(0)\|\leq L\|Y\|$ for all $Y$ with $\|Y\|_2<\delta$. 
Since 
\[|\inpr{\frac{1}{it I-|\Lambda+Y|^2} e_k}{e_i}\inpr{\frac{1}{it I-|\Lambda+Y|^2}(I-Q(Y)) e_j}{e_l}\psi_1(it)|
\leq \frac{|\psi_1(it)|}{|t|\sqrt{t^2+c_2^2}}, \]
the Lebesgue convergence theorem implies  
\begin{align*}
	&\lim_{Y\to 0}
	\frac{1}{2\pi i}\int_{\partial R_{b^2,\varepsilon}} \inpr{\frac{1}{\zeta I-|\Lambda+Y|^2} e_k}{e_i}\inpr{\frac{1}{\zeta I-|\Lambda+Y|^2}(I-Q(Y)) e_j}{e_l}	\psi_1(\zeta)d\zeta\\ \nonumber
	&=\delta_{k,i}\delta_{l,j}\frac{1}{2\pi i}\int_{\partial R_{b^2,\varepsilon}}\frac{\psi_1(\zeta)}{(\zeta-|\lambda_i|^2)(\zeta-|\lambda_j|^2)}d\zeta. 
\end{align*}
Thus it suffices to show 
\[\lim_{Y\to 0}
	\frac{1}{2\pi i}\int_{\partial R_{b^2,\varepsilon}} \inpr{\frac{1}{\zeta I-|\Lambda+Y|^2} e_k}{e_i}\inpr{\frac{1}{\zeta I-|\Lambda+Y|^2}Q(Y)e_j}{e_l}	\psi_1(\zeta)d\zeta=0\]
	
We choose an orthonormal basis $\{v_p\}_{p=1}^n$ of $\C^n$ satisfying $|\Lambda+Y|v_p=s_p(\Lambda+Y)v_p$. 
Since 
\begin{align*}
	&\inpr{\frac{1}{\zeta I-|\Lambda+Y|^2} e_k}{e_i}\inpr{\frac{1}{\zeta I-|\Lambda+Y|^2}Q(Y)e_j}{e_l} \\
	&\sum_{p,q}\inpr{e_k}{v_p}\inpr{v_p}{e_i}
	\inpr{Q(Y)e_j}{v_q}\inpr{v_q}{e_l}\frac{1}{\zeta-s_p(\Lambda+Y)^2}\frac{1}{\zeta-s_q(\Lambda+Y)^2},
\end{align*}
where $p$ runs from 1 to $n-m^g_1$ and $q$ runs from $n-m^a_1+1$ to $n-m^g_1$, we have 
\begin{align*}
	 &\frac{1}{2\pi i}\int_{\partial R_{b^2,\varepsilon}} \inpr{\frac{1}{\zeta I-|\Lambda+Y|^2} e_k}{e_i}\inpr{\frac{1}{\zeta I-|\Lambda+Y|^2}Q(Y)e_j}{e_l}	\psi_1(\zeta)d\zeta  \\
	&= \sum_{p,q}\inpr{e_k}{v_p}\inpr{v_p}{e_i}
	\inpr{Q(Y)e_j}{v_q}\inpr{v_q}{e_l}\psi_1^{[1]}(s_p(\Lambda+Y),s_q(\Lambda+Y)).
\end{align*} 
Since 
\[\|Q(Y)e_j\|_2=\|(Q(Y)-Q(0))e_j\|_2\leq L\|Y\|,\] 
and $|\psi_1^{[1]}(s_p(\Lambda+Y),s_q(\Lambda+Y))|\leq \|\psi_1'\|_\infty<\infty$, we get 
\begin{align*}
	&|\frac{1}{2\pi i}\int_{\partial R_{b^2,\varepsilon}} \inpr{\frac{1}{\zeta I-|\Lambda+Y|^2} e_k}{e_i}\inpr{\frac{1}{\zeta I-|\Lambda+Y|^2}Q(Y)e_j}{e_l}	\psi_1(\zeta)d\zeta|  \\
	&\leq  n^2L \|\psi_1'\|_\infty \|Y\|_2,
\end{align*} 
which finishes the proof. 
\end{proof}

\end{document}
